Die Arbeit $W = m\,g\,h$ hängt nur von der Last und der Höhe ab, nicht von der Zeit. Sie bleibt also \result{\text{gleich}}.

Der Hebevorgang dauert nur noch $t_2 = \left(1-\rtO\right)\cdot t_1$. Die Leistung $P = W/t$ ist umgekehrt proportional zur Zeit:
\begin{align}
 \frac{P_2}{P_1} &= \frac{t_1}{t_2} = \frac{1}{1-\rtO}
\end{align}
Die Leistung nimmt also um
\begin{align}
 \frac{P_2-P_1}{P_1} &= \frac{1}{1-\rtO}-1 = \zP \approx \result{\zPP{0}{2}}
\end{align}
zu, also um etwas mehr als \rtO.
